\documentclass[10pt,a4paper]{amsart}
\usepackage[margin=19mm]{geometry}
\usepackage{amsmath,amssymb,mathtools}
\usepackage[scr=rsfs,cal=euler]{mathalfa}
\usepackage{xfrac}
\usepackage[unicode,hidelinks,bookmarks=false]{hyperref}
\newcommand{\ie}{i.e.\ }
\newcommand{\cf}{cf.\ }
\newcommand{\resp}{respectively\ }
\newcommand{\Subsection}{\subsection}
\providecommand{\nobreakdash}{\nobreak\hbox{-}}
\setlength{\emergencystretch}{2em}
\multlinegap0pt
\usepackage{bm}
\usepackage{bbm}
\usepackage{dsfont}
\usepackage{xspace}
\usepackage{cancel}
\usepackage{mathtools}
\usepackage{amsthm}
\makeatletter
\@ifundefined{theorem}{\newtheorem{theorem}{Theorem}}{}
\@ifundefined{lemma}{\newtheorem{lemma}[theorem]{Lemma}}{}
\makeatother
\title[An Inductive Approach to Basepoint-Freeness of Adjoint Serie]{An Inductive Approach to Basepoint-Freeness of Adjoint Series on Irregular Varietiess}
\author{Houari Benammar Ammar}
\date{Source: arXiv:2408.04733v2 (CC BY 4.0)}
\begin{document}
\maketitle

\noindent\emph{Source and licence.} An Inductive Approach to Basepoint-Freeness of Adjoint Series on Irregular Varietiess. Version arXiv:2408.04733v2 (CC BY 4.0), distributed under the Creative Commons Attribution 4.0 International licence, as recorded in the arXiv licence metadata for this exact version and shown on the abstract page https://arxiv.org/abs/2408.04733v2. Full rights notice: licenses/arxiv-2408.04733-CC-BY-4.0.txt.\par\smallskip

\noindent\textit{Editorial scope.} Counted unit: the Lemma whose source label is \texttt{lemma1} in the article (source0.tex lines 376--378, source label \texttt{lemma1}), delivered together with the article's own complete original proof of it (source0.tex lines 379--423, one \texttt{proof} environment of 45 lines). No mathematics is written, repaired or re-derived: statement and proof are the article's own text line for line. Required context retained uncounted: none. Every definition, display and label of those lines is retained unchanged. Omitted from this package and not redistributed: 1--375, 424--665. Nothing inside a retained excerpt is omitted or altered: every retained body is delivered exactly as the article's lines are written, so no omission receipt of any kind accompanies this package. Citations: the retained excerpts carry no citation command at all, so no bibliography entry of the article is transcribed and no reference is redistributed. Reference adaptation: none was needed and none was made; every reference inside the delivered unit resolves inside it, so no unresolved reference or citation is produced. Printed slips kept literal: no printed word, symbol, hypothesis or number of the counted statement or of its proof was altered. Layout and class: the article's own class is \texttt{sn-jnl} with packages including graphicx, multirow, amsmath, amssymb, amsfonts, amsthm, lmodern, mathrsfs, appendix, xcolor, textcomp, manyfoot, booktabs, algorithm; the wrapper is the lane's pinned \texttt{amsart} editorial wrapper at 10pt on A4, which restates the theorem-like environments the retained text needs and adds the article's own macro definitions that the retained text needs (none), with extra wrapper packages (bm, bbm, dsfont, xspace, cancel, mathtools). The delivered numbering is the unit's own. Assets: no retained span contains a figure environment, a graphics-inclusion command, a PDF-page inclusion, a table or a TikZ picture; no image file is copied into this package, the package declares no asset dependency and no image file is redistributed with it. Licence: version arXiv:2408.04733v2 of this article is distributed under the Creative Commons Attribution 4.0 International licence, as recorded in the arXiv licence metadata for this exact version and shown on the abstract page https://arxiv.org/abs/2408.04733v2. A full rights notice is delivered at licenses/arxiv-2408.04733-CC-BY-4.0.txt. Characters: every retained excerpt is pure ASCII, so the delivered engine, XeLaTeX with its default Latin Modern text font, reports no missing character.
\begin{lemma}\label{lemma1}
Let $X$ be an irregular variety of dimension $n\geq 2$ with $h:X \to Y$ be a morphism to an Abelian variety $Y$. Let $X \xrightarrow{f} Z \xrightarrow{u} Y$ be the Stein factorization of $h$, and let $F$ be a general fiber of $f$. Let $N$ be a divisor on $X$. If $|N_{|_{F}}|$ is basepoint-free and $f_{*}\mathcal{O}_{X}(N)$ is continuously globally generated, then $|N + h^{*}p|$ has no basepoints supported on $F$ for general $p \in \operatorname{Pic}^{0}(Y)$.
\end{lemma}

\begin{proof} Indeed, take a point $x\in F$ such that $z =f(x)$. By assumption, $f_{*}\mathcal{O}_{X}(N)$ is continuously globally generated, that is
for every nonempty open subset $U \subset \operatorname{Pic}^{0}(Z)$, the following direct sum of evaluation maps 
\begin{equation}\label{se1*}
\bigoplus_{p \in U} H^{0}(Z, f_{*}\mathcal{O}_{X}(N)\otimes p) \otimes p^{\vee} \to f_{*}\mathcal{O}_{X}(N)_{|_{z}}
\end{equation}
is surjective. Since $\operatorname{Pic}^{0}(Y) \hookrightarrow \operatorname{Pic}^{0}(Z)$, we have in particular that for every nonempty open subset $U \subset \operatorname{Pic}^{0}(Y)$, the direct sum of evaluation maps
\begin{equation}\label{se1**}
\bigoplus_{p \in U} H^{0}(Z, f_{*}\mathcal{O}_{X}(N)\otimes u^{*}p) \otimes u^{*}p^{\vee} \to f_{*}\mathcal{O}_{X}(N)_{|_{z}}
\end{equation}
is surjective. Furthermore,  \begin{equation}\label{se1}
f_{*}\mathcal{O}_{X}(N)_{|_{z}} = H^{0}(F, \mathcal{O}_{F}(N)) 
\end{equation}
and 
\begin{equation}\label{se2}
H^{0}(Z, f_{*}\mathcal{O}_{X}(N)\otimes u^{*}p) \simeq H^{0}(X, \mathcal{O}_{X}(N) \otimes h^{*}p).
\end{equation}
Combining (\ref{se2}) with (\ref{se1}) and (\ref{se1**}), we obtain that the following sum of maps
\begin{equation}\label{se3}
\bigoplus_{p \in U} H^{0}(X, \mathcal{O}_{X}(N) \otimes h^{*}p) \otimes h^{*}p^{\vee} \to H^{0}(F, \mathcal{O}_{F}(N))
\end{equation}
is surjective. By hypothesis, $|N_{|_{F}}|$ is basepoint-free. Together with (\ref{se3}), this implies that for every nonempty open subset $U \subset \operatorname{Pic}^{0}(Y)$, the following direct sum of evaluation maps
\begin{equation}\label{se4}
\bigoplus_{p \in U} H^{0}(X, \mathcal{O}_{X}(N) \otimes h^{*}p) \otimes h^{*}p^{\vee} \to \mathcal{O}_{X}(N)_{|_{x}}    
\end{equation}
is surjective. In particular, for some $p_{U} \in U$, the map 
\begin{equation}\label{se5}
 H^{0}(X, \mathcal{O}_{X}(N) \otimes h^{*}p_{U}) \otimes h^{*}p_{U}^{\vee} \to \mathcal{O}_{X}(N)_{|_{x}}    
\end{equation}
is surjective. Thus, twisting $(\ref{se5})$ by $h^{*}p_{U}$, we deduce that 
\begin{equation}\label{se6}
 H^{0}(X, \mathcal{O}_{X}(N) \otimes h^{*}p_{U}) \to \mathcal{O}_{X}(N)_{|_{x}}    
\end{equation}
is surjective. Hence, the linear system $$|N + h^{*}p_{U}|$$ has no  basepoints supported on $F$. Now, we define the following subset $S \subseteq \operatorname{Pic}^{0}(Y)$ by 
\[
S := \left\{ p \in \operatorname{Pic}^{0}(Y) \;\middle|\; |N + h^{*}p| \text{ has no basepoints supported on } F \right\}.\]
We claim that \( S \) is dense. Indeed, for any nonempty open subset \( U \subseteq \operatorname{Pic}^{0}(Y) \), one can find an element \( p_{U} \in U \) such that \( |N + h^{*}p_{U}| \) has no basepoints supported on \( F \), as proved. Therefore, \( S \cap U \neq \emptyset \) for every nonempty open subset 
\( U \subseteq \operatorname{Pic}^{0}(Y) \). Hence, \( S \) is dense. 
Otherwise, \( \left(\operatorname{Pic}^{0}(Y) \setminus \overline{S}\right) \cap S \neq \emptyset \), 
a contradiction. It follows that 
\[
|N + h^{*}p|
\]
has no basepoints supported on \( F \) for general 
\( p \in \operatorname{Pic}^{0}(Y) \).
\end{proof}

\end{document}
